\section{Алгоритмы комплексной навигации}
\label{sec:python-complex-nav}

\begin{missionbox}
Глава продолжает миссии Python №\,64--72 (измерения и устойчивость) и №\,91--99 (наведение по метке): ограничение скорости, мёртвая зона, потеря метки и снижение только над целью --- это и есть практическая часть регулятора, о котором идёт речь ниже.
\end{missionbox}

\subsection{Проблема шума и фильтрация}
Данные, получаемые от системы компьютерного зрения, всегда содержат шум. Координаты дрона <<дрожат>> даже если он висит неподвижно. Это происходит из-за дискретности матрицы камеры, изменений освещения и вибраций корпуса.

Если передавать эти <<сырые>> координаты напрямую в полетный контроллер, дрон будет дергаться, пытаясь отработать каждое ложное смещение.

\subsubsection{Фильтр Калмана (Kalman Filter)}
Это стандарт де-факто в робототехнике~\cite{kalman1960}. Фильтр Калмана не просто усредняет значения, он предсказывает, где дрон \textit{должен} оказаться в следующий момент времени, и корректирует это предсказание на основе новых измерений.

Алгоритм работает циклически в два этапа:
\begin{enumerate}
    \item \textbf{Prediction (Предсказание):} На основе физической модели (скорость, ускорение) мы оцениваем новую позицию. При этом растет неопределенность (мы не уверены в прогнозе).
    \item \textbf{Update (Коррекция):} Мы получаем измерение от камеры. Сравниваем его с прогнозом. Чем точнее датчик (меньше шум измерения $R$), тем больше мы верим ему. Чем точнее наша модель (меньше шум процесса $Q$), тем больше верим прогнозу.
\end{enumerate}

\begin{codefiletitle}[python]{Пример 1D фильтра Калмана}{python/advanced_navigation/kalman_1d_example.py}
\end{codefiletitle}

Выше приведен только ключевой метод \texttt{update()} с простейшей моделью «величина не меняется». Полный запускаемый скрипт с классом \texttt{KalmanFilter1D}, экспоненциальным сглаживанием для сравнения и модельным опытом находится в файле \texttt{advanced\_navigation/kalman\_filter.py} --- именно его нужно запускать в заданиях ниже. Он работает на компьютере (нужен \verb|numpy|, для графика --- \verb|matplotlib|): дрон висит на высоте 1,5~м, датчик шумит, а на 50-м шаге дрон «толкнули» на 2,0~м. Программа печатает разброс оценок до скачка и через сколько шагов каждый фильтр «догнал» новую высоту.

\begin{taskbox}[Сравнение фильтров]
\emph{Компьютер.} Запустите скрипт \texttt{kalman\_filter.py}.
\begin{enumerate}
    \item Измените \texttt{MEASUREMENT\_NOISE} (шум измерения $R$) с 0.2 на 0.01, затем на 5.0. Как меняются разброс оценки и время реакции?
    \item Сравните фильтр Калмана и экспоненциальное сглаживание при резком скачке высоты. Кто быстрее реагирует? Подберите \texttt{ALPHA} так, чтобы разброс сглаживания стал как у Калмана.
\end{enumerate}
\expected{с параметрами по умолчанию разброс оценок около 0,03~м против 0,18~м у измерений, реакция на скачок --- 23 шага у Калмана и 18 у сглаживания. При малом $R$ фильтр почти повторяет измерения (быстро, но шумно), при большом --- сглаживает сильнее и реагирует медленнее. Вывод: фильтр со статической моделью всегда меняет плавность на запаздывание.}
\end{taskbox}

\begin{taskbox}[Сглаживаем ошибку наведения]
\emph{Симулятор.} В сцене миссии Python №\,95 (дрон над меткой 12 в точке $(0.6, 0.3, 3)$) в 20 свежих кадрах вычисляйте ошибку $e_x$ и пропускайте её через экспоненциальное сглаживание с $\alpha = 0.3$. Печатайте сырую и сглаженную ошибку. Затем отлетите к $Y = 5$ (метка пропадает, как в миссии №\,97) и вернитесь.
\expected{пока метка видна, сглаженная ошибка меняется плавнее сырой; при потере метки программа \emph{сбрасывает} фильтр (\texttt{value = None}) и обнуляет скорость, а после возвращения начинает сглаживание заново, а не со старого значения.}
\mission{Python}{95--97}
\end{taskbox}

\subsection{Системы координат}
Одна из самых частых ошибок новичков — путаница в системах координат.

\begin{itemize}
    \item \textbf{Body Frame (Связанная система):} оси привязаны к корпусу. В робототехнике чаще всего X смотрит вперёд (нос дрона), Y — влево, Z — вверх. Если дрон поворачивается, оси поворачиваются вместе с ним. Внимание: \verb|pioneer_sdk| в командах \verb|..._body_fixed| понимает связанные оси иначе --- \emph{x вправо, y вперёд}, z вверх.
    \item \textbf{World Frame / NED (North-East-Down):} Ось X смотрит на Север, Y — на Восток, Z — Вниз (к земле). Эта система неподвижна. Локальная система Пионера, в которой мы задавали точки (раздел~\ref{sec:missions}), тоже неподвижна, но устроена иначе: X вперёд от старта, Y влево, Z \emph{вверх}. Если дрон на старте смотрит на север, при переходе между ними меняются знаки Y и Z.
    \item \textbf{Camera Frame (система камеры OpenCV):} X вправо по кадру, Y вниз по кадру, Z --- вперёд из объектива. В ней выражены \texttt{tvec} и \texttt{rvec}. Для нижней камеры ось Z камеры смотрит вниз, и связь с осями дрона задаётся тем, как камера закреплена на корпусе.
    \item \textbf{Marker Frame:} Ось Z обычно смотрит перпендикулярно плоскости маркера (вверх от него).
\end{itemize}

Для навигации нам нужно переводить положение маркера, измеренное камерой, в глобальные координаты (World Frame). Сначала вектор из системы камеры поворачивают в систему корпуса ($R_{cam}$ --- как закреплена камера), затем --- в мировую ($R_{drone}$ --- текущая ориентация дрона):
$$ Pos_{world} = Pos_{drone} + R_{drone} \cdot R_{cam} \cdot Pos_{marker\_camera} $$

\begin{taskbox}[Оси кадра и оси дрона]
\emph{Симулятор.} Повторите опыт миссий Python №\,75 и №\,92 для обеих осей: зависните над меткой 12 на $Z = 3$, запишите ошибку $(e_x, e_y)$, затем одну секунду лететь со скоростью 0,2~м/с сначала по $+X$, потом (после возврата) по $+Y$ (\verb|set_manual_speed|, команда каждые 0,1~с, затем явная остановка) и снова измерьте ошибку.
\expected{таблица «движение дрона --- изменение $e_x$, $e_y$» для нижней камеры симулятора; по ней выбраны знаки поправок, и они совпадают с контроллером \texttt{align} миссии №\,96 ($v_x$ со знаком $+e_x$, $v_y$ со знаком $-e_y$).}
\mission{Python}{75, 92, 96}
\end{taskbox}

\begin{taskbox}[Мировые и связанные оси]
\emph{Симулятор.} Взлетите на $Z = 1.5$. При курсе 0 запомните положение, две секунды подавайте \verb|set_manual_speed(0.3, 0, 0, 0)| (каждые 0,1~с), остановитесь и напечатайте смещение по $X$ и $Y$ (\verb|get_local_position_lps|). Повторите то же с \verb|set_manual_speed_body_fixed(0, 0.3, 0, 0)|. Затем поверните дрон на курс $\pi/2$ (\verb|go_to_local_point(x, y, z, yaw=math.pi/2)|) и проделайте оба опыта ещё раз.
\expected{таблица из четырёх смещений длиной около 0,6~м: у \texttt{set\_manual\_speed} направление смещения при обоих курсах одно и то же (мировая ось $X$), у \texttt{set\_manual\_speed\_body\_fixed} оно поворачивается на 90° вместе с курсом --- дрон летит туда, куда смотрит нос.}
\mission{Python}{45}
\end{taskbox}

\subsection{Точная посадка с PID-регулятором}
Для автоматической посадки на зарядную станцию используется PID-регулятор.
\begin{itemize}
    \item \textbf{P (Proportional):} Чем дальше мы от центра, тем сильнее наклоняемся к нему.
    \item \textbf{I (Integral):} Если нас сносит ветром, мы накапливаем ошибку и наклоняемся сильнее, чтобы компенсировать снос.
    \item \textbf{D (Derivative):} Если мы летим к центру слишком быстро, мы <<тормозим>> заранее, чтобы не пролететь мимо.
\end{itemize}
Команда скорости на каждом шаге:
$$ v = K_p e + K_i \sum e \,\Delta t + K_d \frac{e - e_{\text{пред}}}{\Delta t}, $$
после чего $v$ обязательно ограничивают сверху (миссия №\,70), а внутри мёртвой зоны обнуляют (миссия №\,95).

Программа ниже --- П-регулятор из миссий №\,94--99, собранный в одну посадку: на каждой ступени высоты (3, 2 и 1,5~м) она добивается трёх кадров подряд с меткой в допуске 8 пикселей и только потом снижается дальше; при потере метки скорость обнуляется. Запускайте её в редакторе симулятора в сцене миссии №\,99.

\begin{codefiletitle}[python]{Посадка на метку ступенями (симулятор)}{python/advanced_navigation/staged_landing_sim.py}
\end{codefiletitle}

\begin{warnbox}[При переносе на настоящий дрон]
\verb|Camera(view="bottom")| --- расширение симулятора. На дроне камеру выбирают под его оборудование, а знаки и коэффициенты проверяют заново коротким импульсом (как в миссии №\,92): от того, как закреплена камера, зависит, в какую сторону лететь при $e_x > 0$.
\end{warnbox}

\begin{taskbox}[Посадка в круг]
\emph{Симулятор.} Возьмите \verb|staged_landing_sim.py| в сцене миссии Python №\,99.
\begin{enumerate}
    \item Запустите как есть и запишите точку посадки.
    \item Добавьте D-составляющую: храните прошлую ошибку и прибавляйте $K_d (e - e_{\text{пред}})$ с $K_d = 0.5$ (ошибку берите в долях кадра, как $K_p$). Сравните число кадров до «По центру» на высоте 3~м.
    \item Уменьшите допуск \verb|TOL_PX| до 4 пикселей. Что происходит со временем наведения?
\end{enumerate}
\expected{дрон садится не дальше 0,15~м от метки (в этой сцене она лежит около точки $(0.6; 0.3)$ --- сравните с напечатанной точкой посадки); с D-составляющей наведение заканчивается за меньшее или равное число кадров без раскачки; при слишком узком допуске растёт время, а если метку не удаётся удержать в допуске за 120 попыток, программа печатает «Прервано» и садится.}
\mission{Python}{94--99}
\end{taskbox}

Посадку на незнакомой карте (миссия №\,107) вы уже выполняли в главе~\ref{sec:python-final}: программа, которая берёт положение метки только из кадров, садится на неё где бы та ни стояла, а программа с запомненными координатами --- нет.

\begin{questionbox}[Контрольные вопросы]
\begin{enumerate}
    \item Зачем нужен фильтр Калмана, если можно просто взять среднее арифметическое за 10 кадров?
    \item Что произойдет с PID-регулятором, если перепутать знак у коэффициента P?
    \item В какой системе координат (Body, World или камеры) камера выдает вектор \texttt{tvec}?
    \item Почему при потере метки нельзя продолжать полёт со старой скоростью?
\end{enumerate}
\end{questionbox}

\subsection*{Итоги}
\begin{itemize}
    \item Данные зрения шумят: сырое измерение сглаживают (экспоненциальное сглаживание, фильтр Калмана), а при потере цели фильтр сбрасывают.
    \item Перед расчётом поправки выясняют систему координат: оси кадра, оси корпуса и оси сцены различаются, а \texttt{set\_manual\_speed} и \texttt{set\_manual\_speed\_body\_fixed} понимают оси по-разному.
    \item Регулятор: P --- пропорционально ошибке, I --- накопленная ошибка, D --- скорость её изменения; результат всегда ограничивают, а внутри мёртвой зоны обнуляют.
    \item Точная посадка идёт ступенями: центрирование, подтверждение несколькими кадрами, снижение, снова центрирование.
\end{itemize}
